% !TeX root = closed-systems-from-comparison-completeness.tex
% ============================================================
\chapter*{Orientation to Part II}
\phantomsection
\label{part:intro-quotient}
% ============================================================
\begin{chapterorientation}
\orientationitem{Settled}{Part~I supplies comparison completeness, canonical factorization, and diagonal redundancy.}
\orientationitem{Task}{Part~II determines the quotient semantics forced by that structure and states the resulting closure theorem.}
\orientationitem{Handoff}{The output is the physical quotient and the precise obstruction to primitive subsystem attribution.}
\end{chapterorientation}

Part~II turns the product and diagonal-redundancy results into semantic
structure.  A closed comparison theory cannot treat representative
coordinates, subsystem labels, or supported descriptions as physical
before quotient descent.  The physical object is the quotient-level
content that survives the intrinsic diagonal symmetry.

\Cref{ch:rotation} proves the descent statement and records the twin
representative descriptions that no coherent report can distinguish.
\Cref{ch:closure} then consolidates the preceding argument into the
closure theorem.  At the end of the part, closure is no longer a
guiding phrase; it is the profile-pairing and quotient-semantic package
that later transport constructions must respect.
